\section{Esperança e Momentos}

A esperança de uma variável aleatória representa uma medida de centralidade dos valores
dessa variável.
Intuitivamente, ela é obtida ponderando-se os valores possíveis da variável por suas
respectivas probabilidades, sumarizando a sua imagem a um único valor central.
Na linguagem de medida, a definição formal da esperança é dada a seguir.

\begin{def:esperanca}
Seja $X$ uma variável aleatória e $(\Omega, \mathcal{F}, P)$ espaço de probabilidade.
Defini-se esperança de $X$, ou valor esperado, denotada por $\ev{X}$, como
\[
	\ev{X} = \int_\Omega X\, dP
\]
\end{def:esperanca}

Podemos obter o valor esperado de uma transformação de uma variável aleatória ponderando
os valores da transformação pelas próprias probabilidades de suas respectivas pré-imagens,
em vez da probabilidade dos valores em si, o que é útil quando não conhecemos em detalhes
a distribuição de probabilidade da variável obtida pela transformação.

\begin{thm:esperanca de transformacao}
Seja $X$ variável aleatória e $g: \real \to \real$ uma transformação mensurável.
A esperança da variável $g(X)$ é dada por
\[
	\ev{g(X)} = \begin{cases}
		\int_{-\infty}^\infty g(x) f(x)\, dx, & \text{ se } $X$ \text{ é contínua}  \\
		\sum_{x \in \Omega_X} g(x) P(X = x),  & \text{ se }  $X$ \text{ é discreta} \\
	\end{cases}
\]
\begin{proof}
	Considere o espaço de medida $(\Omega, \mathcal{F}, P)$, tal que
	$g(X): \Omega \to \real$ e
	\[
		\ev{g(X)} = \int_\Omega g(X)\, dP
	\]
	mas pelo Teorema \ref{thm:mudanca de variavel}, uma vez que $g(X) = g \circ X$, então
	\[
		\ev{g(X)} = \int_\real g(x)\, dP_x
	\]
	Se $X$ for contínua, então $f(x) = \frac{d}{d\lambda}P_X$ é densidade de $X$ e fazendo
	a substituição
	$dP_X = f(x) d\lambda$, obtemos
	\[
		\ev{g(X)} = \int_\real g(x) f(x)\, d\lambda = \int_{-\infty}^\infty g(x)f(x)\, dx
	\]
	Por outro lado, se $X$ for discreta, então, pela Proposição \ref{prop:existencia fmp},
	$p_X(x) = \frac{d}{d\mu}P_X$ é a função de massa de $X$, onde $\mu$ é medida de
	contagem, e fazendo a substituição $dP_X = p_Xd\mu$, obtemos
	\[
    	\ev{g(X)} = \int_\real g(x) p_X(x)\, d\mu = \int_{\real \setminus \Omega_X} g(x)p_X(x)\, d\mu +
            \int_{\Omega_X} g(x)p_X(x)\, d\mu
	\]
	onde $\Omega_X$ é o suporte de $X$.
	Uma vez que,
	\[
	    \forall x \in \real,\, ( x \notin \Omega_X \rightarrow p_X(x) = 0 )
	\]
	então
	\[
    	\ev{g(X)} = \int_{\Omega_X} g(x)p_X(x)\, d\mu
	\]
	e, pelo Teorema \ref{thm:integracao contagem},
	\[
    	\ev{g(X)} = \sum_{x \in \Omega_X} g(x)p_X(x)
	\]
\end{proof}
\end{thm:esperanca de transformacao}

O Teorema a seguir enuncia quatro propriedades da esperança enquanto um operador:
linearidade, positividade, monotonicidade e preservação de limites.
Esse resultado serve como base de cálculos utilizando-se o operador de valor esperado.

\begin{thm:propriedades esperanca}
\label{thm:propriedades esperanca}
Seja $X$ variável aleatória, e $g_1$ e $g_2$ transformações mensuráveis.
O operador de esperança satisfaz:
\begin{enumerate}
	\item $\ev{ag_1(X) + bg_2(X)} = a \ev{g_1(X)} + b \ev{g_2(X)}$
	\item $\forall x \in \Omega_X, (g_1(x) \ge 0) \rightarrow \ev{g_1(X)} \ge 0$
	\item $\forall x \in \Omega_X,\ (g_1(x) \ge g_2(x)) \rightarrow \ev{g_1(X)} \ge \ev{g_2(X)}$
	\item $\forall x \in \Omega_X, (a \le g_1(x) \le b) \rightarrow a \le \ev{g_1(X)} \le b$
\end{enumerate}
\begin{proof}
	A propriedade 1 segue da linearidade dos operadores de integral, no caso contínuo, e
	de somatório no caso discreto.
	A propriedade 2 advém da nulidade para o contínuo e  do fechamento dos reais não
	negativos sob adição nos discretos.
	As propriedades 3 e 4 vêm da monotonicidade das integrais e da propriedade aditiva das
	desigualdades.
\end{proof}
\end{thm:propriedades esperanca}

Os momentos de uma variável aleatória medem diferentes aspectos da distribuição dela.
A intuição advém da mecânica, em que os momentos de massa descrevem como a massa
de um sistema está distribuída entorno de um ponto.
Analogamente, os momentos descrevem como a probabilidade dos valores da variável estão
distribuídos.

\begin{def:momento}
Seja $n \in \nat$ e $X$ variável aleatória com esperança $\mu$.
O enésimo momento de $X$, denotado $\mm_(\mu)'$, é definido como
\[
	\mm_n(\mu)' = \ev{X^n}
\]
O enésimo momento central, denotado $\mm_n(\mu)$, é dado por
\[
	\mm_n(\mu) = \ev{(X - \mu)^n}
\]
onde $\mu = \ev{X}$.
\end{def:momento}

Cada valor de $n$ determina um momento de uma variável, de forma que os dois primeiros
possuem interpretações intuitivas:
\begin{itemize}
\item O primeiro momento é o ``centro de massa'' da distribuição, consistindo no próprio
valor esperado da variável.

\item O segundo é a ``inércia'' da variável, medindo o quão disprso estão os seus valores.
Valores altos recebem pesos maiores do que os mais próximos da origem, tal modo que
um pequeno segundo momento diz que as realizações estão conentradas na origem, enquanto
um alto segundo momento diz que eles estão dispersos entorno dela.
\end{itemize}

Os momentos centrais descrevem a distribuição da variável entorno de seu centro de massa,
ou seja, do seu valor esperado:

\begin{itemize}
\item O segundo momento central mede o quão dispersos os valores estão em relação ao valor
esperado.

\item O terceiro momento, por preservar o sinal, descreve a direção em que os valores
se concentram em relação ao valores esperado: caso seja à esquerda, então ele será
negativo; se à direita, então positivo; e caso a distribuição seja aproximadamente
simétrica, então ele estará próximo de zero.

\item O quarto momento central explica a distribuição de valores extremos, pois esses
recebem pesos muito maiores do que valores mais próximos da esperança.
\end{itemize}

Na análise estatística, os principais momentos centrais usados para explicar uma variável
aleatória são o segundo, terceiro e quarto, os quais podem ser normalizados -- com exceção
do segundo -- para criarem medidas descritivas.

\begin{def:momentos importantes}
Se $X$ uma variável aleatória com esperança $\mu$, então estão definidos:
\begin{enumerate}
	\item Variância: $\var{X} = \mm_2(\mu)$
	\item Desvio-padrão: $\dep{X} = \sqrt{\var{X}}$
	\item Assimetria: $\frac{\mm_3(\mu)}{\dep{X}^3}$
	\item Curtose: $\frac{\mm_4(\mu)}{\dep{X}^4} - 3   $
\end{enumerate}
\end{def:momentos importantes}

A variância -- e, por consequência, o desvio-padrão -- é rotineiramente usado com mais
frequência junto da esperança para descrever distribuições.
Assim como o caso da esperança, nos é útil conhecer algumas propriedades da variância
enquanto um operador para fins de cálculos.

\begin{thm:propriedades variancia}
\label{prop:propriedades variancia}
Seja $X$ variável aleatória com esperança $\mu$ e $a \in \real$.
Vale que
\begin{enumerate}
	\item $\var X = \ev X^2 - \mu^2$
	\item $\var{aX} = a^2 \var{X}$
	\item $\var{X + a} = \var{X}$
\end{enumerate}
\begin{proof}
	Com efeito, do Teorema \ref{thm:propriedades esperanca} segue que
	\begin{align*}
		\var{X} & = \ev{(X - \mu)^2}               \\
		        & = \ev{X^2 - 2X\mu + \mu^2}       \\
		        & = \ev{X^2} - 2\mu \ev{X} + \mu^2 \\
		        & = \ev{X^2} - 2 \mu^2 + \mu^2     \\
		        & = \ev X^2 - \mu^2
	\end{align*}
	provando 1.
	Usando isso para 2,
	\begin{align*}
		\var{aX} & = \ev{(aX)^2} - \ev{aX}^2             \\
		         & = \ev{a^2X^2} - \ev{aX} \cdot \ev{aX} \\
		         & = a^2 \ev{X^2} - a^2 \ev{X}^2         \\
		         & = a^2 \var{X}
	\end{align*}
	E a aquela ainda pode ser usada para obter 3:
	\begin{align*}
		\var{X + a} & = \ev{(X + a)^2} - \ev{X + a}^2                           \\
		            & = \ev{X^2} + 2a \ev{X} + a^2 - (\ev{X} + a)(\ev X + a)    \\
		            & = \ev{X^2} + 2a \ev{X} + a^2 - \ev{X}^2 - 2a \ev{X} - a^2 \\
		            & = \ev{X^2} - \ev{X}^2                                     \\
		            & = \var{X}
	\end{align*}
\end{proof}
\end{thm:propriedades variancia}

Uma transformação aplicável a variáveis aleatórias é aquela que subtrai a mede e divide
o resultado pelo desvio-padrão, conhecida como \textbf{padronização}.
O destaque para ela justifica-se pelo fato de que a variável resultante sempre terá
valor esperado igual a $0$ e desvio-padrão igual a $1$, o que é provado na Proposição
\ref{prop:padronizacao}.
Esse resultado é especialmente útil quando temos um conjunto de várias variáveis e queremos
tratá-las na mesma escala.

\begin{prop:padronizacao}
\label{prop:padronizacao}
Se $X$ é uma variável aleatória com esperança $\mu$ e variância $\sigma^2$, então
a variável definida por
\[
    Z = \frac{X - \mu}{\sigma}
\]
possui esperança 0 e variância $1$.
\end{prop:padronizacao}
\begin{proof}
	Com efeito,
	\[
	    \ev{Z} = \frac{1}{\sigma} \cdot \ev{X} - \frac{\mu}{\sigma} = 0
	\]
	enquanto que,
	\[
	    \var{Z} = \ev{\frac{(X - \mu)^2}{\sigma^2}} = \frac{\sigma^2}{\sigma^2} = 1
	\]
\end{proof}

As funções geradoras de momentos (mgf), como seu nome sugere, são usadas para obter os
momentos de uma variável aleatória.
No entanto, na prática, é mais razoável utilizar a própria definição de momento para fins
de computação.
A sua real utilidade está em oferecer uma forma de representar todos os momentos de uma
variável enquanto uma função com poderosas propriedades teóricas.

\begin{def:mgf}
A função geradora de momentos (mgf) da variável $X$ é
\begin{gather*}
	M_X: \nat \to \real \\
	M_X(t) = \ev{e^{tX}}
\end{gather*}
caso o valor esperado exista para $t$ numa vizinhança de $0$.
Caso contrário, dizemos que a mgf não existe.
\end{def:mgf}

A presença de uma exponencial no argumento do valor esperado na definição da mgf pode
ser confusa a primeira vista.
A razão disso é revalada na demonstração do Teorema \ref{thm:derivada do momento}, o qual
enuncia como obter os momentos por meio da mgf: avaliando a drivada de ordem $n$ do
gerador de momentos em torno de $t = 0$.

\begin{thm:derivada do momento}
\label{thm:derivada do momento}
Se $X$ uma variável aleatória cuja mgf existe, então
\[
    \ev{X^n} = \left. \frac{d^n}{dt^n} M_X(t)\right|_{t=0}.
\]
\end{thm:derivada do momento}
\begin{proof}
Se $f(x)$ é a função de densidade de $X$, então
\[
    \left. \frac{d^n}{dt^n} M_X(t)\right|_{t=0} =
    \left. \frac{d^n}{dt^n} \int_{-\infty}^\infty e^{tx} f(x)\, dx \right|_{t=0}
\]
Usando a Regra de Leibniz, que permite diferenciar sobre o sinal de integral,
\[
    \left. \frac{d^n}{dt^n} \int_{-\infty}^\infty e^{tx} f(x)\, dx \right|_{t=0} =
    \int_{-\infty}^\infty \left. \frac{\partial^n}{\partial t^n} e^{tx} f(x)\, dx \right|_{t=0}.
\]
Uma vez que
\[
    \frac{\partial^n}{\partial t^n} e^{tx} = x^n e^{tx},
\]
então
\[
    \left. \frac{d^n}{dt^n} M_X(t)\right|_{t=0} =
    \left. \int_{-\infty}^\infty  x^n e^{tx} f(x)\, dx \right|_{t=0} =
    \int_{-\infty}^\infty x^n f(x)\, dx =
    \ev {X^n}
\]
\end{proof}

O próximo resultado diz como obter a mgf de uma variável obtida por transformação afim
de uma outra.

\begin{prop:mgf da transformacao afim}
Se $X$ e $Y$ é uma variável aleatória tal que $Y = a X + b$, com $a$ e $b$ reais, então
\[
    M_Y(t) = e^{bt} M_X(at)
\]
\end{prop:mgf da transformacao afim}
\begin{proof}
Com efeito,
\begin{align*}
	M_Y(t) &= E[e^{ty}] \\
	&= \ev{e^{(at)x} \cdot e^{tb}} \\
	&= e^{tb} \cdot \ev{e^{(at)x}} \\
	&= e^{tb} \cdot M_X(at)
\end{align*}
\end{proof}

Apesar de serem usados para descrever distribuições de variáveis aleatórias, momentos
não necessariamente as definem.
Em outras palavras, embora duas variáveis possuam a mesma sequência de momentos, não
necessariamente elas possuem a mesma distribuição.
O Teorema \ref{thm:condicao de igualdade de distribuicao e momento} estabelece algumas
condições para que sequências de moementos idênticos permitem dizer que estamos tratando
das mesmas distribuições.

\begin{thm:condicao de igualdade de distribuicao e momento}
\label{thm:condicao de igualdade de distribuicao e momento}
Sejam $X$ e $Y$ variáveis aleatórias com distribuições $F_X(x)$ e $F_Y(y)$ cujos
momentos existem. Valem as seguintes afirmações:
\begin{enumerate}
\item Se os suportes das variáveis são limitados, $F_X(u) = F_Y(u)$ para todo $u \in \real$
se, e somente se, $\ev{X^r} = \ev{Y^r}$ para todo inteiro $r$ não negativo.

\item Se a funções geradoras de momentos das variáveis existem e $M_X(t) = M_Y(t)$ para
todo $t$ em uma vizinhança de $0$, então $F_X(u) = F_Y(u)$ para todo $u \in \real$.
\end{enumerate}
\end{thm:condicao de igualdade de distribuicao e momento}
\begin{proof}
Sejam $f_x$ e $f_Y$ as densidades de $X$ e $Y$ respectivamente.
Utilizando-se a integração por partes para expressar $\ev{X^n}$,
\[
    \ev{X^n} = \int_{-\infty}^\infty u^n f_X(u)\,du =
    u^n F_X(u) - n\int_{-\infty}^\infty u^{n-1} F_X(u)\,du
\]
mas supondo que $F_X(u) = F_Y(u)$, então
\[
    \ev{X^n} = u^n F_Y(u) - n\int_{-\infty}^\infty u^{n-1} F_Y(u)\,du =
    \int_{-\infty}^\infty u^n f_Y(u)\,du =
    \ev{Y^n}
\]
\end{proof}

Aqui é que notamos a utilidade das funções geradoras de momentos: apesar de uma sequência
de momentos não ser suficiente para caracterizar uma distribuição, a convergência de
geradores de momentos permite afirmar a convergência de uma distribuição.

\begin{thm:convergencia mgf}
\label{thm:convergencia mgf}
Seja $(X_i)_{i=1}^n$ uma sequência de variáveis aleatórias com geradores de momentos
$M_{X_i}(t)$ respectivamente.
Se para todo $t$ numa vizinhança de $0$ é satisfeito que
\[
    \lim_{i \to \infty} M_{X_i}(t) = M_X(t)
\]
então há uma única distribuição $F_X(x)$ tal que
\[
    \lim_{i \to \infty} F_{X_i}(x) = F_X(x)
\]
\end{thm:convergencia mgf}

A demonstração desse Teorema está fora do escopo deste texto, pois ela baseia-se no fato
de que $M_X(t)$ é Transformada de Laplace de uma única função de densidade $f(x)$.
Também seria necessário mostar que, dada a convergência das transformadas, as densidades
das variáveis envolvidas satisfazem determinadas hipóteses de continuidade e convergência,
as quais permitem dizer que as próprias densidades convergem para a mesma função $f(x)$.
